\chapter{Đi Học Muộn, Nói Sao Cho Duyên}
\chapterintro{Nhắc mà không xé không khí đầu giờ}{Một câu nói đúng có thể giữ kỷ luật mà không làm học sinh bật chế độ chống đối.}

\section{Em đến đúng lúc lớp đang nhớ em đấy}
\begin{manthoai}{Màn thoại}
\textbf{Thầy/Cô:} Em đến đúng lúc lớp đang nhớ em đấy.\\
\textbf{Học sinh:} Dạ em xin lỗi ạ.\\
\textbf{Thầy/Cô:} Rồi, vào chỗ nhanh để lớp bớt nhớ tiếp.
\end{manthoai}
\begin{dungkhinào}
Hợp với trường hợp đi muộn nhẹ, không cố tình, khi muốn giữ tiết học không bị ngắt mạnh.
\end{dungkhinào}
\begin{nhipthem}
Sau câu vui phải có yêu cầu rõ: vào chỗ, mở vở, theo bài ngay.
\end{nhipthem}
\ngancach

\section{Đi muộn thì đừng muộn luôn cả sự tập trung}
\begin{manthoai}{Màn thoại}
\textbf{Thầy/Cô:} Đi muộn thì đừng muộn luôn cả sự tập trung nhé.\\
\textbf{Học sinh:} Dạ.
\end{manthoai}
\begin{dungkhinào}
Dùng khi học sinh vào lớp rồi nhưng vẫn loay hoay, còn làm ảnh hưởng nhịp chung.
\end{dungkhinào}
\begin{nhipthem}
Giọng nên nhẹ và ngắn. Câu này mạnh ở chỗ nhắc trúng việc cần làm ngay.
\end{nhipthem}
\ngancach

\section{Em nợ lớp một lý do ngắn thôi}
\begin{manthoai}{Màn thoại}
\textbf{Thầy/Cô:} Em nợ lớp một lý do ngắn thôi, không cần nguyên bộ phim.\\
\textbf{Học sinh:} Dạ tại xe hư ạ.\\
\textbf{Thầy/Cô:} Rồi, cảm ơn bản trailer. Giờ vào bài.
\end{manthoai}
\begin{dungkhinào}
Hợp với học sinh có xu hướng kể dài để trì hoãn việc quay lại tiết học.
\end{dungkhinào}
\begin{nhipthem}
Nếu đi muộn lặp lại nhiều lần, sau tiết học vẫn cần nói chuyện riêng chứ không chỉ dùng lời dí dỏm.
\end{nhipthem}
\ngancach

\section{Hôm nay em vào lớp kiểu chạy deadline à?}
\begin{manthoai}{Màn thoại}
\textbf{Thầy/Cô:} Hôm nay em vào lớp kiểu chạy deadline à?\\
\textbf{Học sinh:} Dạ em trễ ạ.\\
\textbf{Thầy/Cô:} Rồi, deadline bây giờ là ngồi xuống và theo bài ngay.
\end{manthoai}
\begin{dungkhinào}
Hợp với học sinh vào lớp vội nhưng không có dấu hiệu chống đối, cần kéo ngay vào việc chính.
\end{dungkhinào}
\begin{nhipthem}
Giữ giọng gọn, tránh để lớp cuốn theo màn đối đáp dài ngay đầu giờ.
\end{nhipthem}
\ngancach

\section{Trễ rồi, mình đừng trễ luôn phần học nữa nhé}
\begin{manthoai}{Màn thoại}
\textbf{Thầy/Cô:} Trễ rồi, mình đừng trễ luôn phần học nữa nhé.\\
\textbf{Học sinh:} Dạ.
\end{manthoai}
\begin{dungkhinào}
Dùng khi học sinh đã vào lớp nhưng còn chậm chạp mở đồ, khiến tiết bị kéo giãn thêm.
\end{dungkhinào}
\begin{nhipthem}
Câu này ngắn, rõ, không mỉa mai quá mức, rất hợp để nhắc nhanh rồi tiếp tục dạy.
\end{nhipthem}
